% Section 07: Guidance Artifact
% Answers sRQ4: How should effective GenAI practices be systematized into guidance?
% Mode: Prescriptive (guidance) — distinct from descriptive mode in section 06

\section{Guidance Artifact}\label{sec:artifact}

\subsection*{Introduction}

This section presents the guidance artifact answering sRQ4: \textit{``How should
effective GenAI practices for RA development be systematized into guidance for EA
Practitioners?''} It turns the Delphi findings from
\textit{Results and Analysis} (Section~\ref{sec:results}) into prescriptive guidance that Enterprise
Architecture (EA) practitioners can apply when using Generative AI (GenAI) in
Reference Architecture (RA) development. The consensus-variability profiles
from \textit{Results and Analysis} (Section~\ref{sec:results}) directly inform where the artifact can offer
directive recommendations and where guidance must remain adaptive. The artifact
is organized by practice category (per the panel's preference), followed by
cross-cutting guidance, governance adaptation, and a maturity model. The
following subsections detail the artifact:

\localtableofcontents\clearpage

%% ================================================================
%% 7.1 Introduction and Scope
%% ================================================================
\subsection{Introduction and Scope}\label{subsec:artifact-scope}

The guidance artifact presented here is the final version, incorporating Round~3
priority ratings (Table~\ref{tab:guidance-priorities}), artifact critique
feedback (Section~\ref{subsubsec:artifact-critique}), and heuristic refinement
data (Table~\ref{tab:heuristic-importance}). The intended audience is Enterprise
Architecture (EA) practitioners broadly defined: enterprise architects, solution
architects, platform engineers, data architects, and senior engineers who
participate in RA development. This scope matches the panel diversity described
in Section~\ref{subsec:panel-overview}.

The artifact is organized by practice category rather than by RA lifecycle phase
or experience level, consistent with the panel's explicit preference
(Section~\ref{subsubsec:artifact-critique}). One participant stated: ``By
practice seems fine, because if you do by lifecycle, then you will repeat
yourself'' (R3 Brainstorm). Lifecycle phase mapping received the lowest priority
rating among all guidance elements ($M = 2.43$,
Table~\ref{tab:guidance-priorities}), consistent with this choice.

Each practice category carries a guidance level (\textit{directive},
\textit{contextual}, or \textit{adaptive}) based on panel agreement and
variability patterns from Table~\ref{tab:practice-importance}. Categories where
the panel agreed closely ($SD < 1.0$) receive directive guidance; those with
moderate variability ($1.0 \leq SD \leq 1.2$) receive contextual guidance; and
those with wide disagreement ($SD > 1.2$) receive adaptive guidance. Two
overrides apply: Adversarial Review ($SD = 0.97$) is contextual rather than
directive because its value depends on domain expertise, and Deep Workflow
Integration ($SD = 1.17$) is adaptive because its low mean and abstention
pattern point to adoption barriers that vary across settings. Practice
descriptions ($M = 4.71$) and verification patterns ($M = 4.14$), the two
highest-priority guidance elements, structure each practice block.

Each element carries a provenance label. \textit{Delphi consensus} marks
elements the panel voted on or explicitly validated. \textit{Researcher
synthesis} marks researcher interpretation anchored in Delphi data.
\textit{Researcher proposal} marks informed suggestions not validated by the
panel. Where an element draws on both sources, both are indicated. External
literature corroborating each practice is examined separately in
Section~\ref{subsec:disc-triangulation}.

%% ================================================================
%% 7.2 Practice Catalog
%% ================================================================
\subsection{Practice Catalog}\label{subsec:practice-catalog}

Table~\ref{tab:practice-catalog} summarizes the seven validated practice
categories with their consensus profiles, guidance levels, and primary
verification patterns. The following subsections detail each practice.

\begin{table}[!ht]
    \centering
    \caption{Practice Catalog Summary}\label{tab:practice-catalog}
    \begin{tabular}{|p{4.2cm}|c|c|p{2.2cm}|p{4cm}|}
        \hline
        \textbf{Practice Category} & \textbf{$M$} & \textbf{$SD$} & \textbf{Guidance Level} & \textbf{Primary Verification Pattern} \\
        \hline
        System / Codebase Analysis & 4.00 & 0.82 & Directive & Cross-validate against source systems \\
        \hline
        Research \& Exploration & 3.60 & 1.17 & Contextual & Verify claims against primary sources \\
        \hline
        Adversarial Review / Stress Testing & 3.60 & 0.97 & Contextual & Domain expertise threshold check \\
        \hline
        Governance / Compliance & 3.50 & 1.27 & Adaptive & Line-by-line verification in compliance contexts \\
        \hline
        First Draft \& Documentation & 3.10 & 1.10 & Contextual & Cross-check references and facts \\
        \hline
        Stakeholder Communication & 3.10 & 1.37 & Adaptive & Validate organizational context manually \\
        \hline
        Deep Workflow Integration & 2.89 & 1.17 & Adaptive & Incremental trust building \\
        \hline
    \end{tabular}
\end{table}

Throughout the catalog, practice descriptions and verification patterns carry
Delphi consensus; experience indicators are researcher proposals.

\textbf{System/Codebase Analysis} \textit{[Directive].}
This practice involves using GenAI to analyze existing systems, codebases, and
architectural artifacts to build understanding of the current state. Practitioners
should apply this practice during Phase~2 (Current State Analysis) of the RA
lifecycle, particularly for legacy system comprehension, dependency mapping, and
cross-repository analysis. System/Codebase Analysis received the highest
importance rating with the lowest variability
(Table~\ref{tab:practice-importance}), indicating strong consensus that this is
the most valuable GenAI application in RA work.

\textit{Verification pattern:} Cross-validate GenAI-generated system
descriptions against source systems directly. As the Delphi panel established,
all participants delegated information gathering to GenAI while retaining
judgment (Section~\ref{subsubsec:practice-discovery}). For codebase analysis
specifically, practitioners should verify that referenced components, services,
and dependencies exist in the actual system. One Round~3 participant reported
that GenAI ``pointed to a folder that simply didn't exist''
(Section~\ref{subsubsec:trust-calibration}), demonstrating why physical
constraint checking is essential.

\textit{Experience indicators:} Novice practitioners should start with
well-documented codebases where verification is straightforward. Advanced
practitioners extend this practice to legacy systems and cross-repository
analysis where GenAI's compression value is highest.

\textbf{Research \& Exploration} \textit{[Contextual].}
This practice involves using GenAI as a research assistant for exploring design
alternatives, technology options, and architectural trade-offs during target
state design. Seven of 19 Round~1 practices mapped to Phase~4 (Target State
Design) in research-assistant mode
(Section~\ref{subsubsec:practice-discovery}).

\textit{Verification pattern:} Verify claims, statistics, and technology
recommendations against primary sources. The cross-validation heuristic received
the highest importance rating across all heuristics ($M = 4.57$,
Table~\ref{tab:heuristic-importance}). P2 reported verifying ``60--70\% of
specific claims'' against primary sources (R1 Workbook). Practitioners should
treat GenAI research outputs as hypotheses requiring confirmation, not as
authoritative findings.

\textit{Experience indicators:} Start with well-scoped questions verifiable
against established sources; progress to multi-model workflows and iterative
refinement.

\textbf{Adversarial Review / Stress Testing} \textit{[Contextual].}
This practice involves deliberately using GenAI to challenge, critique, and
stress-test architectural decisions. The Delphi panel validated contradiction
seeding as a verification heuristic ($M = 3.33$,
Table~\ref{tab:heuristic-importance}), and domain expertise threshold rated
$M = 3.86$. Adversarial Review requires the practitioner to actively elicit
critique rather than accept GenAI's default collaborative stance.

Adversarial Review encompasses a family of techniques that practitioner discourse
calls \textit{Architecture Stress Testing}. Table~\ref{tab:stress-testing}
presents seven named sub-practices identified through researcher synthesis. The
Delphi panel validated contradiction seeding and domain expertise threshold as
techniques; the remaining techniques were identified from practitioner discourse
and are corroborated in Section~\ref{subsec:disc-triangulation}.

\begin{table}[!ht]
    \centering
    \caption{Architecture Stress Testing Techniques}\label{tab:stress-testing}
    \begin{tabular}{|p{4cm}|p{9cm}|}
        \hline
        \textbf{Technique} & \textbf{Description} \\
        \hline
        AI Sparring Partner & Iterative adversarial dialogue where GenAI challenges design rationale \\
        \hline
        Devil's Advocate Prompting & Instructing GenAI to argue against the proposed architecture \\
        \hline
        Decision Stress Testing & Sequential methodology: identify blind spots, steelman alternatives, conduct pre-mortem \\
        \hline
        AI Pre-Mortem & ``The architecture already failed; why?'' Prospective hindsight analysis \\
        \hline
        Skeptical Reviewer & Role-based persona critique simulating a senior architect review \\
        \hline
        Architectural Gap Analysis & Structural completeness analysis against standards and patterns \\
        \hline
        Epistemic Stress Testing & Challenges confidence levels behind decisions, not the decisions themselves \\
        \hline
    \end{tabular}
\end{table}

\textit{Verification pattern:} Adversarial Review requires domain expertise as a
prerequisite. The panel confirmed that ``a lot of the hallucination issue can
only be identified if you know things are wrong from experience already''
(Section~\ref{subsubsec:trust-calibration}). Practitioners should evaluate
GenAI-generated critiques against their own domain knowledge, recognizing that
GenAI lacks the ability to model strategic adversarial behaviour across
organizational boundaries.

\textit{Experience indicators:} Novice practitioners should start with
structured techniques such as AI Pre-Mortem and Skeptical Reviewer, which
provide clear prompting templates. Advanced practitioners develop intuition for
when GenAI critique reflects genuine design weaknesses versus pattern-matching
artifacts.

\textbf{Governance/Compliance Checking} \textit{[Adaptive].}
This practice involves using GenAI to verify architectural artifacts against
governance standards, compliance requirements, and organizational policies.
Governance/Compliance emerged during live Round~2 discussion as the seventh
category (Section~\ref{subsubsec:practice-validation}), with high variability
($SD = 1.27$) reflecting context-dependent value across organizational settings.
The adaptive guidance level acknowledges that governance structures vary
substantially across organizations.

\textit{Verification pattern:} In compliance contexts, verify every detail
line-by-line. The Delphi panel rated this heuristic $M = 4.20$ with the lowest
variability among Round~2 heuristics. GenAI can scan documents against regulatory
requirements, but practitioners must verify completeness and accuracy manually.
Compliance errors carry disproportionate organizational risk.

\textit{Experience indicators:} Begin with low-risk compliance contexts;
progress to automated governance pipelines where organizational maturity
permits.

\textbf{First Draft \& Documentation} \textit{[Contextual].}
This practice involves using GenAI to produce initial drafts of architectural
documentation (solution designs, architecture decision records (ADRs),
technical specifications, and assessment reports) for subsequent human editing
and validation. Six Round~1 practices employed the First Draft Generator
interaction mode (Section~\ref{subsubsec:practice-discovery}). The Compression
Not Generation framing (Section~\ref{subsubsec:cross-cutting}) applies directly:
GenAI compresses existing information into structured documents rather than
creating novel architectural knowledge.

\textit{Verification pattern:} Cross-check all references, links, and factual
claims in generated documentation. The panel's experience included hallucinated
benchmarks and non-existent references that nearly entered planning documents
(Section~\ref{subsubsec:cross-cutting}). Practitioners should treat generated
drafts as starting points requiring architectural judgment, not finished
artifacts.

\textit{Experience indicators:} Start with well-structured document types
(meeting minutes, status reports) before progressing to architecture decision
records and metaprompt libraries.

A related emerging practice, Spec-Driven Development, extends this
documentation-first principle; its external evidence base is examined in
Section~\ref{subsec:disc-triangulation}.

\textbf{Stakeholder Communication} \textit{[Adaptive].}
This practice involves using GenAI to translate architectural content for
different audiences (executive summaries, team-specific views, non-technical
explanations) and to synthesize stakeholder inputs into structured architectural
artifacts. Stakeholder Communication exhibited the highest variability of all
categories ($SD = 1.37$, Table~\ref{tab:practice-importance}), indicating
substantial disagreement about GenAI's value for this purpose.

\textit{Verification pattern:} Validate organizational context manually before
communicating GenAI-assisted outputs. The panel established that organizational
context (team dynamics, political constraints, historical decisions) remains
exclusively human-retained knowledge
(Section~\ref{subsubsec:cross-cutting}). GenAI-generated stakeholder
communications risk being ``technically logical but organisationally
nonsensical'' (P3, R1 Workbook) if organizational context is not carefully
injected.

\textit{Experience indicators:} Limit initial use to summarization tasks
(meeting minutes, document distillation) where source material constrains
outputs; extend to audience-specific prompting and RAG-based approaches as
familiarity grows.

\textbf{Deep Workflow Integration} \textit{[Adaptive].}
This practice involves embedding GenAI into the practitioner's daily workflow
through IDE integration, agent context files, custom hooks, Model Context
Protocol (MCP) servers, and multi-agent orchestration. Deep Workflow Integration
received the lowest mean rating ($M = 2.89$, one abstention,
Table~\ref{tab:practice-importance}), consistent with its dependence on
individual maturity and organizational tooling policies.

The low rating likely reflects adoption barriers (cognitive load of tooling
setup, organizational resistance, and infrastructure prerequisites) rather than
low practice value. Two panel members operated in Cyborg mode with deep tooling
integration (Section~\ref{subsubsec:practice-discovery}), showing that
practitioners who overcome these barriers achieve sustained workflow benefits.

\textit{Verification pattern:} Build trust incrementally. The incremental trust
building heuristic rated $M = 4.29$ (Table~\ref{tab:heuristic-importance}),
making it the second-highest heuristic and particularly relevant for deep
integration where automated workflows reduce human oversight touchpoints.
Practitioners should progressively expand GenAI autonomy from low-stakes tasks
to higher-stakes ones as calibration develops.

\textit{Experience indicators:} Novice practitioners should begin with
single-tool integration (IDE assistant) before progressing to custom context
files and agent orchestration. Advanced practitioners design multi-agent
architectures with specialized agents for different architectural tasks.

%% ================================================================
%% 7.3 Cross-Cutting Guidance
%% ================================================================
\subsection{Cross-Cutting Guidance}\label{subsec:cross-cutting-guidance}

The practice catalog addresses individual practice categories. Three
cross-cutting concerns apply across all seven categories: trust calibration,
context engineering, and skills atrophy mitigation. These concerns emerged from
the cross-cutting patterns identified in Section~\ref{subsubsec:cross-cutting}
and the verification heuristics refined across all three Delphi rounds
(Section~\ref{subsubsec:heuristics}).

\subsubsection{Trust Calibration Framework}\label{subsubsec:trust-framework}

Trust calibration is the foundational cross-cutting concern. The Delphi panel
established a universal verification mandate: all eight Round~1 participants
described explicit verification practices, with no participant accepting GenAI
outputs without review (Section~\ref{subsubsec:cross-cutting}). The
cross-validation heuristic received the highest importance rating of any
heuristic ($M = 4.57$, $SD = 0.53$, Table~\ref{tab:heuristic-importance}).

Table~\ref{tab:trust-spectrum} presents a trust spectrum synthesized from Delphi
data and external evidence (Researcher synthesis). The panel validated the
underlying verification heuristics, but did not vote on the specific trust
levels or their ordering. The spectrum structure is a researcher
interpretation anchored by the panel's verification experiences and the
organizational context boundary established in
Section~\ref{subsubsec:cross-cutting}. This spectrum translates the
verification heuristics from Section~\ref{subsubsec:heuristics} into task-level
trust guidance, connecting the empirical sense-making patterns (sRQ3) to the
systematized guidance framework (sRQ4).

\begin{table}[!ht]
    \centering
    \caption{Trust Spectrum for GenAI in RA Development}\label{tab:trust-spectrum}
    \begin{tabular}{|p{2.5cm}|p{5.5cm}|p{5.5cm}|}
        \hline
        \textbf{Trust Level} & \textbf{Task Type} & \textbf{Verification Approach} \\
        \hline
        Highest & Information compression: summaries, boilerplate, dependency graphs & Spot-check representative samples \\
        \hline
        High & Documentation drafts, concept explanation, test generation & Review structure and key claims \\
        \hline
        Medium & Code refactoring, common patterns, meeting synthesis & Verify against organizational standards \\
        \hline
        Medium-Low & Design trade-off exploration, technology recommendations & Cross-validate with primary sources \\
        \hline
        Low & Architecture decisions, novel business logic & Independent domain expert review \\
        \hline
        Very Low & Security-critical code, compliance assessment & Line-by-line verification \\
        \hline
        Lowest & Organizational judgment: politics, team dynamics, legal & Retain entirely; do not delegate \\
        \hline
    \end{tabular}
\end{table}

Trust calibration is bidirectional. The Delphi panel identified instances where
GenAI outperformed manual approaches; dependency graph calculations ``were
actually better than our manual diagrams'' (R3 Brainstorm,
Section~\ref{subsubsec:trust-calibration}). Practitioners should recalibrate
trust upward when evidence warrants, not only guard against errors. The ``junior
colleague'' mental model recurred across the panel: GenAI outputs merit the same
supervisory attention as work from a capable but inexperienced team member.

Two prerequisites anchor the framework. First, trust calibration is teachable but
requires domain expertise as a foundation
(Section~\ref{subsubsec:teachability}). Practitioners ``had to build a deep
understanding before applying it'' (R3 Brainstorm). Second, calibration improves
through the frontier learning pattern: six of eight participants discovered
GenAI's limitations through direct experience with errors
(Section~\ref{subsubsec:cross-cutting}). These findings align with
\textcite{dellacqua_navigating_2023}, who demonstrated that knowledge workers
operating on GenAI's ``jagged technological frontier'' must learn where AI
capability boundaries fall through practice.

\subsubsection{Context Engineering for EA}\label{subsubsec:context-engineering}

\textcite{rajasekaran_context_engineering_2025} defines context engineering as
``the set of strategies for curating and maintaining the optimal set of tokens
during LLM inference.'' The Delphi panel validated the importance of this
practice through the context management heuristic ($M = 4.14$, $SD = 0.90$,
Table~\ref{tab:heuristic-importance}), the third-highest rated heuristic.

Context engineering narrows the distance between what practitioners know and
what GenAI can access. Seven of eight participants identified organizational
context (team dynamics, political constraints, historical decisions) as
outside GenAI capability (Section~\ref{subsubsec:cross-cutting}). Effective
context engineering does not eliminate this boundary but narrows it by making
more organizational knowledge machine-accessible.

For EA practitioners, context engineering manifests in three areas (the panel
validated context management as important; the three-part structure is
researcher synthesis). First,
\textit{agent context files}: project-level configuration files such as
CLAUDE.md and AGENTS.md that encode architectural standards, design principles,
and organizational constraints
\parencite{chatlatanagulchai_agent_readmes_2025}. An empirical study of 2,303
such files found that architecture content appeared in 67.7\% of repositories,
with files undergoing multiple commits and evolving every one to three days.
These files function as a new category of architectural artifact: both
human-readable documentation and machine-executable governance.

Second, \textit{strategic information ordering} mitigates the ``lost in the
middle'' effect, where language models attend less to information in the middle
of long contexts \parencite{liu_lost_middle_2024}. Practitioners should place
critical architectural constraints at the beginning and end of context, with
supporting detail in the middle. Modular rules files (separate context documents
for different architectural concerns) keep individual contexts focused and
within effective processing ranges.

Third, \textit{context scope management} addresses what to include and what to
exclude. Effective context includes: architectural principles and constraints,
relevant standards and patterns, system-specific technical context, and
decision history (ADRs). Effective context excludes: organizational politics
and interpersonal dynamics (which GenAI cannot reliably process), exhaustive
historical documentation (which dilutes signal), and ambiguous or contradictory
requirements (which amplify hallucination risk).

\subsubsection{Anti-Atrophy Strategies}\label{subsubsec:anti-atrophy}

Practitioners risk losing skills they stop exercising. The skills GenAI
automates (cold-start writing, exhaustive manual search, cover-to-cover
documentation reading) are the same skills through which practitioners built
their expertise
\parencite{kabashkin_cognitive_atrophy_2025}. The Delphi panel rated ``skills
used less'' as the second-most significant practice change ($M = 3.90$,
Table~\ref{tab:change-significance}), confirming that practitioners recognize
this dynamic in their own experience.

Three deliberate friction practices mitigate competency atrophy (the panel
validated the atrophy concern; the three mitigation strategies are researcher
proposals). First,
\textit{AI-free design sessions}: periodically conducting architectural work
without GenAI assistance to maintain foundational skills. This practice is
analogous to how pilots train for manual flight despite autopilot capabilities.
Second, \textit{attempt independently first}: engaging with a problem for 15--30
minutes before consulting GenAI, ensuring the practitioner develops and maintains
their own mental model before seeing an AI-generated one. Third, \textit{pair
programming mindset}: treating GenAI interaction as collaborative work rather
than delegation, maintaining active cognitive engagement throughout. The
practitioner drives architectural thinking while GenAI handles information
retrieval and expression.

These strategies connect to the role identity preservation pattern
(Section~\ref{subsubsec:cross-cutting}). All eight Delphi participants viewed
GenAI as an efficiency tool rather than a replacement for professional judgment.
As P4 stated: ``I did this work for years without GenAI. Efficiency would
decrease but effectiveness would not'' (R1 Workbook). Anti-atrophy practices
operationalize this stance by ensuring that practitioners maintain the
independent capability that underpins effective GenAI supervision.

%% ================================================================
%% 7.4 Governance Adaptation
%% ================================================================
\subsection{Governance Adaptation}\label{subsec:governance-adaptation}

Governance adaptation for GenAI in RA development rests on thinner Delphi
evidence than the practice catalog and cross-cutting guidance. Governance
patterns received the lowest significance ratings among cross-role dynamics
(Table~\ref{tab:cross-role}), and Governance/Compliance Checking exhibited the
second-highest variability ($SD = 1.27$,
Table~\ref{tab:practice-importance}). We present three governance themes that
emerged from Delphi data triangulated with external evidence, while
acknowledging that most organizations have not yet formally adapted their
governance structures.

\textbf{Gate-to-guardrail transition} \textit{(Researcher synthesis).} Traditional EA governance operates through
periodic architectural review boards (ARBs) that evaluate decisions at defined
stage gates. GenAI's acceleration of artifact production (practitioners producing
drafts in minutes rather than days) challenges this periodic model. The Delphi
panel's low rating for inside-out governance ($M = 2.00$,
Table~\ref{tab:cross-role}) suggests that most organizations have not yet
confronted this shift formally. \textcite{ettinger_ea_dynamic_capability_2025}
argues that EA frameworks have significant limitations addressing GenAI
requirements, noting the tension between innovation enablement and compliance
maintenance. In our view, organizations should consider transitioning from
periodic gate reviews to continuous guardrail enforcement, where architectural
constraints are encoded and checked automatically rather than evaluated
episodically.

\textbf{Architecture-as-Code enforcement} \textit{(Researcher proposal).} Encoding architectural constraints as
machine-readable rules (TypeScript validators, fitness functions, CI/CD
checks) enables continuous enforcement without manual review bottlenecks. This
approach bridges the Governance/Compliance practice category with Deep Workflow
Integration by embedding governance directly into development workflows. When
architectural rules are machine-readable, GenAI agents can read and apply them
before generating code, reducing violation rates. Early implementations
demonstrate executable ADRs where architectural decisions become automated
rules in build pipelines \parencite{ford_evolutionary_architectures_2023, zdun_sustainable_decisions_2013}. This represents a shift from
documenting governance intent to operationalizing it.

\textbf{AI disclosure and provenance} \textit{(Researcher synthesis).} As GenAI-generated artifacts become
prevalent, governance frameworks need mechanisms for identifying which
components were AI-generated or AI-augmented. The ``Ghost Architect''
dynamic (where ``everybody can generate an architecture now,''
Section~\ref{subsubsec:cross-role}) creates traceability challenges. AI
disclosure policies flag AI-generated components for calibrated review,
enabling review boards to apply appropriate scrutiny. Provenance tracking
becomes particularly important for compliance-sensitive domains where the
origin of architectural decisions carries regulatory implications.

%% ================================================================
%% 7.5 Maturity Model
%% ================================================================
\subsection{Maturity Model}\label{subsec:maturity-model}

The Delphi data revealed that practitioners operate at different levels of GenAI
integration, from ad-hoc experimentation to deep workflow embedding. Round~1
workbooks showed this variation clearly: two participants operated in Cyborg
mode with MCP servers and custom hooks, while others used chat-based interaction
exclusively (Section~\ref{subsubsec:practice-discovery}).
Table~\ref{tab:maturity-model} presents a four-stage maturity model derived from
this observed variation, mapped to the practice catalog and cross-cutting
guidance. The four-stage structure is a researcher proposal; the progression
assumption that practitioners naturally advance from basic to sophisticated
integration has Delphi consensus ($M = 4.71$,
Table~\ref{tab:maturity-progression}).
This structure follows established maturity model conventions.
\textcite{becker_developing_2009} propose a procedure model for
theoretically founded maturity model development, noting that over a
hundred such models have been created across IT management domains.
\textcite{poeppelbuss_what_2011} characterize maturity models as staged
representations of anticipated evolutionary paths, where each stage
builds on the previous one. A recognized critique of such models is
that the assumed sequential progression lacks empirical causal validation
\parencite{poeppelbuss_what_2011, van_steenbergen_maturity_2011};
Section~\ref{subsec:maturity-model} addresses this by grounding the
four stages in observed practitioner variation rather than prescriptive
theory.

\begin{table}[!ht]
    \centering
    \caption{Maturity Model for GenAI in RA Development}\label{tab:maturity-model}
    \begin{tabular}{|c|p{3cm}|p{4.5cm}|p{4.5cm}|}
        \hline
        \textbf{Stage} & \textbf{Label} & \textbf{Characteristics} & \textbf{Observable Indicators} \\
        \hline
        1 & Experimental & Ad-hoc GenAI use for individual tasks; chat-based interaction; no structured verification & Occasional use for research or drafting; no shared practices; trust calibration through trial and error \\
        \hline
        2 & Systematic Personal & Structured personal workflows; context engineering basics; deliberate verification & Personal prompt libraries; agent context files; anti-atrophy practices; consistent trust spectrum application \\
        \hline
        3 & Team-Level Integration & Shared practices across team; team-level agent context files; Architecture-as-Code adoption & Shared context files in version control; team verification standards; collective decision capture workflows \\
        \hline
        4 & Organizational Embedding & Governance adaptation; enterprise context engineering; agentic workflows & Automated architectural guardrails; AI disclosure policies; multi-agent orchestration; enterprise MCP infrastructure \\
        \hline
    \end{tabular}
\end{table}

The SP2 maturity progression vote, collected asynchronously after the Round~3
session (Table~\ref{tab:maturity-progression}), provides partial empirical
grounding. The panel rated natural progression from basic to sophisticated
integration at $M = 4.71$ (variability 22\%), supporting the model's
progression assumption. The proposition that short experience can equal long
experience scored $M = 2.86$, indicating that progression requires sustained
engagement. Per-category measurement scored $M = 3.86$ (variability 49\%),
suggesting practitioners advance at different rates across practice categories.
Context scored the same ($M = 3.86$) but appears to modulate pace rather than
replace experience.

The model is diagnostic, not prescriptive. Practitioners can use it to locate
their current practice and identify growth areas; progression through all four
stages is not necessary for every context. The progression assumption has
panel support, though the four-stage structure itself remains
researcher-proposed. \textit{Data Limitations} (Section~\ref{subsec:data-limitations}) addresses the
methodological caveats, and the \textit{Discussion} (Section~\ref{sec:discussion}) interprets the
implications for EA maturity theory.

